% Threat Model and Trust Assumptions — Article 1
% Formal adversary model (ADV-1..ADV-4). Every attack referenced here maps to a
% row in article1/paper/tables/adversary_attack_mapping.csv and to a raw result
% file under article1/results/raw/. No claim without a raw file.

\section{Threat Model and Trust Assumptions}
\label{sec:threat}

\subsection{System and trust boundary}
We consider a multi-tenant Kubernetes cluster in which some worker nodes are AMD
SEV-SNP Confidential VMs (CVMs). A sensitive AI-inference workload must run only
on a node whose attestation evidence has been \emph{independently verified},
is \emph{fresh}, and is \emph{not revoked}. The trust chain has three roles that
mirror the IETF RATS architecture~\cite{birkholz2023rats}: a node-resident
\emph{attester} (the \texttt{node-attestation-agent}), a central
\emph{verifier} (the \texttt{central-verifier}, the sole writer of
\texttt{AttestationEvidence}), and \emph{relying parties} (the
attestation-aware scheduler and the \texttt{verify-placement} tool).

\paragraph{Trusted computing base (TCB).}
We trust: the Kubernetes API server and etcd; the AMD SEV-SNP hardware root of
trust and the Microsoft Azure Attestation (MAA) verifier within the standard
SEV-SNP threat model; the central-verifier's appraisal logic; and the
scheduler's Ed25519 signing key. The confidentiality/integrity guarantees of
the SEV-SNP CVM itself are assumed per its specification and are supported by
formal analysis of the SEV-SNP software interface~\cite{paradzik2024sevsnp}.

\paragraph{Explicitly out of scope.}
A malicious or compromised Kubernetes API server or etcd; a cloud provider that
violates the SEV-SNP hardware guarantees; micro-architectural side
channels~\cite{li2021cipherleaks}; physical attacks; and compromise of the
scheduler's private signing key. These are orthogonal, hardware- or
platform-level threats addressed by prior work and by the SEV-SNP TCB, not by an
orchestration-layer defense.

\subsection{Adversary classes}
We define four adversaries of increasing capability.

\begin{description}
  \item[ADV-1 — Namespace-scoped application user.] Can create/modify pods and
  policy-adjacent objects within a namespace it controls, and set arbitrary
  pod labels, annotations, \texttt{nodeSelector}, and \texttt{runtimeClassName}.
  It cannot forge attestation evidence.
  \item[ADV-2 — Compromised non-confidential node.] Controls a worker node that
  is \emph{not} a CVM and its node-agent ServiceAccount. It may attempt to
  submit attestation material for itself or to write evidence directly.
  \item[ADV-3 — Compromised confidential node (post-attestation).] Controls a
  genuine CVM after it has attested, e.g. via a runtime compromise. It may try
  to reuse or replay evidence, or to keep serving after its evidence is revoked
  or expires.
  \item[ADV-4 — Object/label manipulation adversary.] Can race mutable cluster
  state: modify a matched \texttt{ConfidentialInferencePolicy}, flip a label,
  revoke evidence, or tamper with a placement token \emph{between} the
  scheduler's Filter and Bind phases (a TOCTOU adversary). It cannot forge
  valid Ed25519 signatures.
\end{description}

\subsection{Security properties}
The design targets six properties, each falsifiable by an attack:
\emph{Placement Safety} (a sensitive pod binds only to a node with valid
evidence), \emph{Evidence Freshness} (stale/revoked evidence is rejected, even
if it becomes stale mid-scheduling), \emph{Identity Binding} (the decision is
bound to the exact pod, node, policy and evidence), \emph{Label Resistance}
(attacker-forgeable metadata cannot grant placement), \emph{Runtime
Consistency} (the effective runtime matches policy, and simulated runtimes are
refused in production), and \emph{Verifiable Placement} (the decision carries an
offline-verifiable cryptographic token).

\subsection{Attack catalog}
Table~\ref{tab:adversary} (source:
\texttt{article1/paper/tables/adversary\_attack\_mapping.csv}) enumerates
attacks A1--A10, the adversary that mounts each, the targeted property, the
expected defense, and the raw-result artifact that evidences the outcome.
The principal security evaluation uses the real AKS SEV-SNP campaign. Local
\texttt{kind} runs are CI/regression artifacts only and are not counted as paper
security results. Attacks not yet re-run on AKS are identified as unit/regression
coverage rather than real-cluster evidence.
Crucially, A10 confirms that a node-agent principal \emph{cannot} write
\texttt{AttestationEvidence}: a forged \texttt{RawAttestationReport} is appraised
by the central verifier to \texttt{evidenceMode=unverified},
\texttt{verificationStatus=Failed}, and never yields real evidence.

\begin{table}[t]
  \centering
  \caption{Adversary/attack mapping (data:
  \texttt{tables/adversary\_attack\_mapping.csv}). Prop.\ = targeted security
  property.}
  \label{tab:adversary}
  \footnotesize
  \begin{tabular}{@{}llp{3.4cm}@{}}
    \toprule
    \textbf{Attack} & \textbf{Adv.} & \textbf{Property / defense} \\
    \midrule
    A1 & ADV-4 & Label Resistance: forged label, no evidence $\to$ Pending \\
    A2 & ADV-1 & Evidence Freshness: stale evidence filtered \\
    A3 & ADV-1 & Placement Safety: revoked evidence filtered \\
    A4 & ADV-1 & Placement Safety: wrong-TEE evidence filtered \\
    A5 & ADV-1 & Runtime Consistency: non-conf.\ runtime denied \\
    A6 & ADV-4 & Runtime Consistency: PreBind policy-hash re-check \\
    A7 & ADV-3 & Evidence Freshness: PreBind revocation re-check \\
    A8 & ADV-4 & Verifiable Placement: tampered token rejected \\
    A9 & ADV-1 & Runtime Consistency: simulated runtime denied in prod \\
    A10 & ADV-2/3 & Placement Safety + Identity Binding: forged evidence $\to$ unverified \\
    \bottomrule
  \end{tabular}
\end{table}

% IP REVIEW REQUIRED BEFORE SUBMISSION
% (Placement token / signed AIPlacementDecision are described in this model.)
